\documentclass{article}

%-------------------------%
% PREAMBUŁA --------------%
%-------------------------%
\usepackage[utf8]{inputenc}
\usepackage[OT4]{polski}
\usepackage{caption}
\usepackage{amsmath}
\usepackage{tikz}
\usepackage{subcaption}
\usepackage{tabularx}
\usepackage{caption}
\usepackage{array}
\usepackage{hyperref}
\usepackage{graphicx}
\usepackage{fancyhdr}
\usepackage[justification=centering]{caption}
\usepackage[margin=60pt]{geometry}
\usepackage{longtable}



%\usepackage[T1]{fontenc}
%\usepackage{amsthm}
%\theoremstyle{plain}
%\usepackage{enumitem}
%\captionsetup[table]{skip=-4pt}
%\usepackage{tabulary}
%\usepackage{sidecap}
%\usepackage{gensymb}
%\usepackage{wrapfig}
%\usepackage{lastpage}
%\usepackage{gensymb}
%\usepackage{lipsum}
%\pagestyle{fancy}


\newlength{\RoundedBoxWidth}
\newsavebox{\GrayRoundedBox}
\newenvironment{GrayBox}[1][\dimexpr\textwidth-4.5ex]%
   {\setlength{\RoundedBoxWidth}{\dimexpr#1}
    \begin{lrbox}{\GrayRoundedBox}
       \begin{minipage}{\RoundedBoxWidth}}%
   {   \end{minipage}
    \end{lrbox}
    \begin{center}
    \begin{tikzpicture}%
       \draw node[draw=black,fill=black!1,rounded corners,%
             inner sep=2ex,text width=\RoundedBoxWidth]%
             {\usebox{\GrayRoundedBox}};
    \end{tikzpicture}
    \end{center}}

%------- Typy kolumn do ustawiania szerokości ------- %
\newcolumntype{L}[1]{>{\raggedright\let\newline\\\arraybackslash\hspace{0pt}}m{#1}}
\newcolumntype{C}[1]{>{\centering\let\newline\\\arraybackslash\hspace{0pt}}m{#1}}
\newcolumntype{R}[1]{>{\raggedleft\let\newline\\\arraybackslash\hspace{0pt}}m{#1}}








\begin{document}

%-------------------------%
%Tabela nagłówkowa -------%
%-------------------------%

\begin{table}[h]
\begin{tabular}{|l|l|l|l|l|l|}
\hline
\begin{tabular}[c]{@{}l@{}}
Wydział:\\ WFiIS\end{tabular} &
\multicolumn{2}{l|}{\begin{tabular}[c]{@{}l@{}}
Imię i nazwisko:\\ 1. Axel Zuziak\\ 2. Marcin Węglarz \end{tabular}} 
& Rok \textbf{II}         
& Grupa \textbf{B}          
& Zespół \textbf{03}      \\ \hline
\textbf{\begin{tabular}[c]{@{}l@{}}
LABOLATORIUM\\ TECHNIK \\ JĄDROWYCH\end{tabular}} &
\multicolumn{4}{l|}{ Temat:\textbf{ Statystyczny charakter rozpadów promieniotwórczych.}
	 }                                                                                                                       & \multicolumn{1}{c|}{\begin{tabular}[c]{@{}c@{}}
Nr ćwiczenia\\ \textbf{1+9}
\end{tabular}} \\ \hline
\begin{tabular}[c]{@{}c@{}}
Data wykonania:\\ 04.03.2015
\end{tabular} &
\begin{tabular}[c]{@{}c@{}}
 Data oddania:\\ 18.03.2015
\end{tabular} &
\begin{tabular}[c]{@{}c@{}}
Zwrot do poprawy: \\ 
\end{tabular}                                 
& 
\begin{tabular}[c]{@{}c@{}}
Data oddania: \\ 
\end{tabular}&
\begin{tabular}[c]{@{}c@{}}
Data zaliczenia:\\
\end{tabular}& 
\begin{tabular}[c]{@{}c@{}}
OCENA:\\
\end{tabular}       
\\ & & & & & \\ \hline
\end{tabular}
\end{table}
%-----Koniec tabeli------

%%%%%%%%%%%% Wstęp %%%%%%%%%%%%
\section{Wstęp teoretyczny}
\section{Aparatura i wykonanie ćwiczenia}
\begin{itemize}
	\item \textbf{Neutronowy miernik wodoru}
	\item \textbf{Wzmacniacz impulsowy}
	\item \textbf{Zasilacz wysokiego napięcia}
	\item \textbf{Analizator amplitudy}
	\item \textbf{Przelicznik}
\end{itemize}

Ćwiczenie rozpoczęto od zapoznania się z aparaturą pomiarową. Przed przystąpieniem do wykonywania pomiarów zmierzono wszystkie wymiary geometryczne zarówno próbek wzorcowych: $W1,W2,...W7$ jak i tych, dla których wyznaczano gęstość wodoru: $P1,P2,...P7$. Zanotowano wagi poszczególnych próbek. Wykonano pomiary przy pomocy neutronowego miernika wodoru dla wszystkich próbek.

\section{Wyniki pomiarów i obliczenia.}

\subsection{Wyznaczenie objętości i gęstości badanych próbek.}
Przy pomocy linijki zmierzono wymiary geometryczne badanych próbek. Wyniki pomiarów i obliczeń przedstawiono w tabeli \ref{wynikiW}.

\begin{table}[h!]
\centering
\label{wynikiW}
\caption{Wymiary geometryczne próbek wzorcowych.}
\begin{tabular}{|C{2cm}|C{2cm}|C{2cm}|C{2cm}|}\hline
	Próbka & Waga [g] & Objętość [cm$^3$] & Zawartość $H$ [\%] \\ \hline
		W7	&	298,4	&	346,40	&	15,1 \\ \hline
		W6	&	436,25	&	361,28	&	8,89 \\ \hline
		W5	&	414,5	&	345,58	&	8,08 \\ \hline
		W4	&	478,44	&	338,70	&	6,42 \\ \hline
		W3	&	538,28	&	346,40	&	5,45 \\ \hline
		W2	&	434,46	&	337,72	&	2,05 \\ \hline
		W1	&	766,3	&	361,28	&	0 \\ \hline

\end{tabular}
\end{table}

\begin{table}[h!]
\centering
\label{wynikiP}
\caption{Wymiary geometryczne badanych próbek}
\begin{tabular}{|C{2cm}|C{2cm}|C{2cm}|C{2cm}|}\hline
	Próbka & Waga [g] & Objętość V [cm$^3$] & U(V) [cm$^3$] \\ \hline
		P1	&	181,31	&	331,89	&	10,14 \\ \hline
		P2	&	371,6	&	345,58	&	10,46  \\ \hline
		P3	&	511,34	&	345,58	&	10,46 \\ \hline
		P4	&	590,2	&	310,37	&	9,78 \\ \hline
		P5	&	593,58	&	338,70	&	10,30 \\ \hline
		P6	&	416,6	&	346,40	&	10,40 \\ \hline
		P7	&	439,47	&	361,28	&	10,67 \\ \hline

\end{tabular}

\end{table}
\newpage
\subsection{Cechowanie neutronowego miernika wodoru}
Do otrzymanych wyników (tabela \ref{cechowanie_tab}) dopasowano prostą metodą regresji korzystając z programu Gnuplot:
\begin{equation}
	J = a_0 + a_1\rho_H
\end{equation}
 Otrzymano wartości:
 \begin{equation*}
 	a_1 = 430112 \pm 1,451\cdot 10^4
 \end{equation*}
 \begin{equation*}
 	a_0 = 9487,86 \pm 1273
 \end{equation*}
Zatem otrzymana krzywa cechowania przyjmuje postać:
\begin{equation*}
	J = 430112\rho_H + 94,87,86
\end{equation*}



\begin{figure}[h!]
	\fontsize{6}{8}\selectfont % zmniejszam czcionke
	\centering
	\resizebox{1.0\textwidth}{!}{\input{cechowanie.tex}}	
	\caption{Krzywa cechowania neutronowego miernika wodoru}
	\label{krzywa}
\end{figure}


%%%%%%%%%%%% BIBLIOGRAFIA %%%%%%%%%%%%
\newpage
\begin{thebibliography}{9}
	
	
	\bibitem{1}
	\url{http://nucleardata.nuclear.lu.se/toi/}

	
	\bibitem{DzK}
	B. Dziunikowski, S.J. Kalita
	\emph{Ćwiczenia laboratoryjne z jądrowych metod pomiarowcyh}, Wydawnictwa AGH, Kraków 1995
	
	\bibitem{przybycien}
	Mariusz Przybycień \emph{Tablice Statystyczne}, dostęp on-line\\
	\url{http://home.agh.edu.pl/~mariuszp/wfiis_stat/tablice_ps_wir.pdf}
\end{thebibliography}
\vspace{2cm}








\end{document}





%\textsc{http://nucleardata.nuclear.lu.se/toi/}
%do bibliografii






